\newlength{\h}
\setlength{\h}{.9cm}

\newcommand{\enc} {\mathsf{AES} }
\newcommand{\drawbc}[5]{%
	% #1: position
	% #2: node name
	% #3: label
	% #4: color
	% #5: size (width and height of the square)

	% Draw square node
	\node[draw, rectangle, minimum width=#5, minimum height=#5,
		inner sep=0pt, align=center, anchor=center, color=#4] (#2) at (#1) {\vphantom{Ay}#3};

	% Triangle perfectly embedded on left side
	\draw[fill=#4, draw=#4]
	($(#2.west) + (0.02, 0.3*#5)$) --  % upper point
	($(#2.west) + (0.02, -0.3*#5)$) -- % lower point
	($(#2.west) + (0.15cm, 0)$) --     % inner tip (fixed width)
	cycle;
}



\newcommand{\BCbloc}[4][]{%
% Dessine le noeud rectangle
\node[#1, draw, rectangle, minimum width=\h, minimum height=\h, anchor=center] (#3) at #2 {#4};
\coordinate (#3/center) at (#3.center);
\coordinate (#3/north) at (#3.north);
\coordinate (#3/east) at (#3.east);
\coordinate (#3/south) at (#3.south);
% Coordonnées du coin inférieur gauche du rectangle pour placer le triangle
\draw[fill=black]
($(#3.north west)!0.15!(#3.north east) - (0, 0.025)$)  -- % Milieu supérieur du côté gauche
($(#3.north east)!0.15!(#3.north west) - (0, 0.025)$) -- % Milieu inférieur du côté gauche
($($(#3.north west)!0.5!(#3.north east)$) - (0, 0.6*0.4*\h)$) -- % Point intérieur (sommet du triangle)
cycle; % Ferme le triangle
}
